\documentclass[11pt,a4paper]{article}
\usepackage[margin=27mm]{geometry}
\usepackage{longtable,booktabs,array,fvextra}
\usepackage{amsmath,amsthm,amssymb,fontspec,unicode-math,hyperref}
\setmainfont{STIXGeneral}
\setmathfont{STIX Two Math}
\setmonofont{Menlo}
\hypersetup{hidelinks}
\usepackage{newunicodechar}
\newfontfamily\numeralfallback{STIX Two Text}
\newunicodechar{∎}{\ensuremath{\blacksquare}}
\newunicodechar{⁰}{{\numeralfallback\char"2070\relax}}
\newunicodechar{₀}{{\numeralfallback\char"2080\relax}}
\newunicodechar{₁}{{\numeralfallback\char"2081\relax}}
\newunicodechar{₂}{{\numeralfallback\char"2082\relax}}
\newunicodechar{₃}{{\numeralfallback\char"2083\relax}}
\newunicodechar{₄}{{\numeralfallback\char"2084\relax}}
\newunicodechar{⁻}{{\numeralfallback\char"207B\relax}}
\newunicodechar{⁸}{{\numeralfallback\char"2078\relax}}
\newunicodechar{ⁱ}{{\numeralfallback\char"2071\relax}}
\newunicodechar{⁴}{{\numeralfallback\char"2074\relax}}
\newunicodechar{⁷}{{\numeralfallback\char"2077\relax}}
\newunicodechar{⁶}{{\numeralfallback\char"2076\relax}}
\newunicodechar{⁵}{{\numeralfallback\char"2075\relax}}
\newunicodechar{₈}{{\numeralfallback\char"2088\relax}}
\newunicodechar{ᵢ}{\textsubscript{i}}
\newunicodechar{ⱼ}{\textsubscript{j}}
\newunicodechar{ₜ}{\textsubscript{t}}
\providecommand{\tightlist}{\setlength{\itemsep}{0pt}\setlength{\parskip}{0pt}}
\newtheorem{theorem}{Theorem}
\newtheorem{proposition}[theorem]{Proposition}
\newtheorem{lemma}[theorem]{Lemma}
\newtheorem{corollary}[theorem]{Corollary}
\theoremstyle{remark}\newtheorem{remark}[theorem]{Remark}
\newcommand{\gen}[1]{\langle #1\rangle}
\newcommand{\normalgen}[1]{\langle\!\langle #1\rangle\!\rangle}
\newcommand{\Z}{\mathbb Z}
\newcommand{\F}{\mathbb F}
\setlength{\emergencystretch}{2em}

\providecommand{\passthrough}[1]{#1}

\title{Translation lattices and marked finite quotients}
\author{Human Compression Project}
\date{}
\begin{document}
\maketitle
\section{Introduction}
Many finite polyhedra in the companion paper come from one of two infinite affine groups. The parameter changes a lattice of translations, while the finite linear group and the distinguished dihedral subgroups stay fixed. Consequently, the proof of the intersection property can be reduced to a finite list of translation differences and then applied to every period at once.

Two issues must be separated. An affine model satisfying a presentation is initially only a quotient of the presented group. We first identify its full translation kernel, using an exact subgroup presentation. We then determine which period quotients preserve the marked subgroup intersections. The compatible-product construction handles a different operation: changing adjacent orders while retaining an already understood marked quotient.

We include complete proofs of the kernel, deletion and square-group statements used in that construction. They are reused from the project's accepted Topic A theory, with their labels retained for traceability; they are paid parts of this independent common paper. No theorem about forest markings, nilpotent enumeration, or arbitrary coset geometries is silently invoked. The affine arguments reorganize the abelian-normal-subgroup method already used by Hou--Feng--Leemans. They replace separate seed calculations by a common point quotient and a single kernel presentation for each affine group. This is a statement about proof reuse, not historical priority.
\section{Rank-three markings}\label{rank-three-markings}

Throughout, \(x^g=g^{-1}xg\) and \([x,y]=x^{-1}y^{-1}xy\).

A rank-three involutory marking is an ordered generating tuple of nonidentity involutions. For a subset
\(I\) of its indices, \(G_I\) denotes the subgroup on those markers. The
intersection property is \(G_I\cap G_J=G_{I\cap J}\). We retain the standard
regular-polytope correspondence. The following standard rank-three criterion is included with its elementary proof.

\textbf{Rank-three test.} For three distinct involutions with commuting
endpoints and adjacent products of order at least three, IP is equivalent
to \(G_{01}\cap G_{12}=G_1\).

\emph{Proof.} Necessity is immediate. For sufficiency, the endpoint subgroup
is \(G_{02}=\{1,r_0,r_2,r_0r_2\}\). If \(G_{02}\cap G_{01}\) contains
\(r_2\) or \(r_0r_2\), then \(r_2\in G_{01}\cap G_{12}=G_1\), contradicting
distinctness. Thus this intersection is \(G_0\); dually
\(G_{02}\cap G_{12}=G_2\). These equalities exclude every generator from
the subgroup on the other two. Intersections involving singletons now
follow, and intersections involving the empty or full index set are
immediate. This covers all subset pairs. \(\square\)

\section{Marked intersections and supported extensions}\label{marked-intersections-and-supported-extensions}

A marked group is a group \(G\) together with a finite indexed generating family \((r_i)_{i\in I}\). Except where stated, the markers need not be involutions or distinct. For \(S\subseteq I\), put
\[
G_S=\langle r_i:i\in S\rangle,\qquad G_\varnothing=1.
\]
We call these subgroups parabolics, retaining their index sets even when two of them happen to coincide. The marking has the intersection property, or IP, if
\[
G_S\cap G_T=G_{S\cap T}\quad(S,T\subseteq I).
\tag{M1}
\]
A marked homomorphism sends each source marker to the corresponding target marker. Its image on \(G_S\) is therefore the target parabolic with index set \(S\).

\textbf{Lemma UA.K0 (lifting one intersection).} Let \(q:H\to Q\) be a homomorphism and let \(C\leq A\cap B\), where \(A,B,C\leq H\). If
\[
q(A)\cap q(B)=q(C),
\]
then
\[
A\cap B=C\quad\Longleftrightarrow\quad A\cap B\cap\ker q\leq C.
\tag{M2}
\]

\emph{Proof.} The forward implication is immediate. For the converse, let \(x\in A\cap B\). Its image belongs to \(q(C)\), so choose \(c\in C\) with \(q(c)=q(x)\). Then \(xc^{-1}\in A\cap B\cap\ker q\), which is contained in \(C\) by hypothesis. Multiplying by \(c\) gives \(x\in C\). The inclusion \(C\leq A\cap B\) gives equality. \(\square\)

\textbf{Theorem UA.K (kernel profiles).} Let \(q:H\twoheadrightarrow Q\) be a marked epimorphism and suppose that the marking of \(Q\) has IP. Write \(K=\ker q\) and \(K_S=K\cap H_S\). Then \(H\) has IP if and only if
\[
K_S\cap K_T=K_{S\cap T}\quad(S,T\subseteq I).
\tag{M3}
\]

\emph{Proof.} Apply UA.K0 to \(A=H_S\), \(B=H_T\), and \(C=H_{S\cap T}\). The image hypothesis is exactly IP in \(Q\). The subgroup on the left of the containment in (M2) is \(K_S\cap K_T\). Its containment in \(H_{S\cap T}\) is equivalent to its containment in \(K_{S\cap T}\); the reverse inclusion always holds. This proves each required intersection and its converse. \(\square\)



A deletion for \(S\subseteq I\) is an endomorphism \(e_S:G\to G\) satisfying \(e_S(r_i)=r_i\) for \(i\in S\) and \(e_S(r_i)=1\) otherwise. Requiring these values matters: a general retraction onto \(G_S\) need not be a deletion.

\textbf{Proposition UA.D (Boolean deletions).} If a marked group has a deletion \(e_S\) for every \(S\subseteq I\), then
\[
e_Se_T=e_{S\cap T},\qquad \operatorname{im}e_S=\operatorname{Fix}(e_S)=G_S.
\tag{D1}
\]
The marking has IP, and
\[
\ker e_S=\langle\!\langle r_i:i\notin S\rangle\!\rangle_G,
\qquad G=(\ker e_S)\rtimes G_S.
\tag{D2}
\]

\emph{Proof.} The two maps in the first equality agree on every generator. The image is generated by precisely the retained markers, and each of those markers is fixed; this proves the image and fixed-point statements. An element of \(G_S\cap G_T\) is fixed by \(e_Se_T=e_{S\cap T}\), so belongs to \(G_{S\cap T}\). The reverse inclusion is immediate.

Let \(M\) be the normal closure in (D2). Clearly \(M\leq\ker e_S\). In \(G/M\), every omitted marker is trivial, so the image of any word equals the image of its deletion. A word in \(\ker e_S\) is therefore trivial in \(G/M\), proving equality. Finally, for \(g\in G\), the factorization \(g=(g e_S(g)^{-1})e_S(g)\) has factors in \(\ker e_S\) and \(G_S\). Their intersection is trivial because \(e_S\) fixes \(G_S\). This is the asserted semidirect product. \(\square\)

We now assemble extensions whose extra kernel is invisible on each proper part of its support. The groups and kernels may be infinite and nonabelian.

\textbf{Theorem UA.SK (supported-kernel realization).} Let \(I\) be finite and let \(\mathcal F\) be a finite family of distinct subsets of \(I\). Let \(B=\langle b_i:i\in I\rangle\) have IP. For every \(S\in\mathcal F\), suppose the following data are given:

\begin{enumerate}
\def\labelenumi{\arabic{enumi}.}
\tightlist
\item
  A deletion homomorphism \(p_S:B\to B_S\), fixing \(b_i\) for \(i\in S\) and killing \(b_i\) otherwise.
\item
  A group \(E_S=\langle e_{S,i}:i\in S\rangle\) and a marked epimorphism \(q_S:E_S\to B_S\).
\item
  Injectivity of \(q_S\) on \((E_S)_T=\langle e_{S,i}:i\in T\rangle\) for every \(T\subsetneq S\).
\end{enumerate}

Put \(N_S=\ker q_S\), and form the group
\[
H=\left\{(b,(x_S)_{S\in\mathcal F})\in
B\times\prod_{S\in\mathcal F}E_S:
q_S(x_S)=p_S(b)\text{ for every }S\right\}.
\tag{K1}
\]
For \(i\in I\), define \(h_i\) to have base coordinate \(b_i\) and \(S\)-coordinate \(e_{S,i}\) when \(i\in S\), and \(1\) otherwise. These elements generate all of \(H\). If \(\pi:H\to B\) is the base projection and \(H_L=\langle h_i:i\in L\rangle\), then, for every \(L\subseteq I\),
\[
\pi(H_L)=B_L,\qquad
\ker\pi=\prod_{S\in\mathcal F}N_S,\qquad
\ker\pi\cap H_L=\prod_{\substack{S\in\mathcal F\\S\subseteq L}}N_S.
\tag{K2}
\]
The factors on the right are embedded as pure coordinate subgroups. The marking of \(H\) has IP. If \(B\) and all \(E_S\) are finite, then
\[
|H_L|=|B_L|\prod_{\substack{S\in\mathcal F\\S\subseteq L}}|N_S|.
\tag{K3}
\]

\emph{Proof.} The compatibility equations in (K1) are preserved by products and inverses, so \(H\) is a subgroup of the displayed direct product. The prescribed tuple \(h_i\) satisfies those equations by the literal deletion rule. Let \(H'\) be the subgroup generated by these tuples. Its base image is \(B\), and the projection of \(H'_S=\langle h_i:i\in S\rangle\) onto \(E_S\) is onto.

We first prove, by induction on inclusion among the finitely many supports, that every pure coordinate element of \(N_S\) belongs to \(H'_S\). Fix \(n\in N_S\). Choose a word on the markers indexed by \(S\) which evaluates to \(n\) in \(E_S\), and evaluate the same word in \(H'\), obtaining \(w\in H'_S\). Its base coordinate lies in \(B_S\) and is \(q_S(n)=1\).

Consider its coordinate at another support \(T\). Because the word uses only markers from \(S\), that coordinate lies in \((E_T)_{S\cap T}\). Its image under \(q_T\) is \(p_T(1)=1\). If \(T\not\subseteq S\), then \(S\cap T\subsetneq T\), and the proper-parabolic injectivity hypothesis forces this coordinate to be \(1\). The only possible nonidentity coordinates of \(w\) are consequently at supports \(T\subseteq S\). Its \(S\)-coordinate is the prescribed \(n\), while every coordinate at a proper support \(T\subsetneq S\) belongs to \(N_T\).

For each such proper support, the induction hypothesis supplies the pure coordinate element having exactly that value in \(N_T\), as an element of \(H'_T\leq H'_S\). Multiply \(w\) by the inverses of those elements. Distinct pure coordinate subgroups commute in the ambient direct product, so these operations remove the unwanted coordinates without changing the \(S\)-coordinate. The result is the pure element \(n\) in coordinate \(S\), and remains in \(H'_S\). A minimal support has no smaller coordinates to remove, so this argument includes the initial step. No commutativity within an individual kernel has been used.

The kernel of \(\pi\) in the full group \(H\) is exactly the direct product of the \(N_S\), by (K1). The induction places that entire kernel in \(H'\). Since \(H'\) already maps onto \(B\), every \(h\in H\) has the same base coordinate as some \(h'\in H'\); then \(h(h')^{-1}\in\ker\pi\leq H'\). Thus \(H'=H\).

The image equality \(\pi(H_L)=B_L\) follows on generators. If \(x\in H_L\cap\ker\pi\), its \(S\)-coordinate lies in \((E_S)_{L\cap S}\) and in \(N_S\). This coordinate is trivial whenever \(S\not\subseteq L\), again by proper-parabolic injectivity. Conversely, when \(S\subseteq L\), the isolation argument placed the entire pure \(N_S\) in \(H_S\leq H_L\). This proves all three assertions in (K2).

The coordinate products in (K2) satisfy
\[
\left(\prod_{S\subseteq L}N_S\right)\cap
\left(\prod_{S\subseteq M}N_S\right)
=\prod_{S\subseteq L\cap M}N_S.
\]
Theorem UA.K now proves IP in \(H\). More explicitly, an element of \(H_L\cap H_M\) has base in \(B_{L\cap M}\); subtracting a lift from \(H_{L\cap M}\) leaves an element in both displayed kernel products, hence in the kernel product for \(L\cap M\). This also shows exactly where IP of the base is used. Finally (K3) is the order formula for the surjection \(H_L\to B_L\) with the kernel just computed. \(\square\)

\textbf{Corollary UA.SK.Local (recognition on an antichain).} Under UA.SK, if \(\mathcal F\) is an antichain under inclusion, then the coordinate projection \(H_S\to E_S\) is a marked isomorphism for every \(S\in\mathcal F\).

\emph{Proof.} It is onto on generators. If an element maps to \(1\), its base coordinate also vanishes, since the base value of a word on \(S\) is the \(q_S\)-image of its \(S\)-coordinate. Formula (K2) says that the remaining kernel in \(H_S\) has only the factor \(N_S\): the antichain excludes every other support contained in \(S\). The \(S\)-coordinate detects that factor, so the element is \(1\). \(\square\)

\textbf{Corollary UA.SK.Delete (derived deletions).} If, in addition, \(B\) has deletions for every subset of \(I\), then every \(E_S\) and \(H\) has all its Boolean deletions.

\emph{Proof.} For \(T\subsetneq S\), the restriction \(q_S:(E_S)_T\to B_T\) is a marked isomorphism: it is onto on generators and injective by hypothesis. Compose \(q_S\), the deletion \(B_S\to B_T\), and the inverse of this isomorphism. This gives a homomorphism \(E_S\to(E_S)_T\) that fixes exactly the \(T\)-markers and kills the others. For \(T=S\) use the identity. Proposition UA.D supplies the Boolean identities.

For \(L\subseteq I\), apply the base deletion \(e_L\) and, in coordinate \(S\), the local deletion onto \(S\cap L\). The two sides of each compatibility equation agree on the local markers; the base deletions commute by UA.D. Thus the coordinatewise operation preserves (K1). On \(h_i\) it fixes \(h_i\) if \(i\in L\), and kills it otherwise. It is the required deletion of \(H\). \(\square\)

\section{Square presentations and their translation periods}\label{square-presentations-and-their-translation-periods}

Let
\[
 a(x,y)=(-x,y),\quad b(x,y)=(y,x),\quad c(x,y)=(x,1-y).
\]
\textbf{Proposition UA.Affine (the complete square presentation).} These generate \(W=\mathbb Z^2\rtimes D_8\), of type \(\{4,4\}\), with
\(ac=ca\). Its presentation is
\[
 W=\langle a,b,c\mid a^2=b^2=c^2=(ab)^4=(bc)^4=(ac)^2=1\rangle. \tag{S1}
\]
\emph{Proof.} Put \(d=bab\), \(y_0=cd\), and
\(x_0=by_0b\). The relations give \([a,d]=1\), \([a,y_0]=1\),
\(dy_0d=y_0^{-1}\), \(ax_0a=x_0^{-1}\), and conjugation by \(b\) swaps
\(x_0,y_0\). Substitution of \(bc=x_0ab\) makes \((bc)^4=1\) equivalent
to \([x_0,y_0]=1\). Thus every word collects to \(x_0^u y_0^v M\),
\(M\in\langle a,b\rangle=D_8\). In the affine model \(x_0,y_0\) are the
unit translations, so these forms are distinct. This proves (S1). \(\square\)

The normal form also fixes our affine multiplication convention: the rightmost transformation acts first. The symbol \(D_8\) denotes the group of signed permutation matrices in two coordinates, of order eight. In all dihedral notation below the subscript denotes the order of the group.

\textbf{Proposition UA.Square (finite square seeds).} For \(n\ge5\), define the invariant lattice
\[
 \Lambda_n=\begin{cases}
  2^e\mathbb Z^2,&n=2e+3,\\
  2^e\langle(1,1),(1,-1)\rangle,&n=2e+4,
 \end{cases}\qquad e\ge1,
\]
and put \(H_n=W/\Lambda_n\). Its order is \(2^n\); the full linear
\(D_8\) coordinate is retained, even if its action on the vertex quotient
is not faithful. Its pair orders are \(4,4,2\). It has all Boolean deletions and IP. Its third lower-central subgroup is \(2\mathbb Z^2/\Lambda_n\), of order \(2^{n-5}\).

\emph{Proof.} The lattice indices are respectively \(2^{2e}\) and \(2^{2e+1}\), so the normal form gives the stated orders. The three linear pair products have orders four, four and two; hence the marked pair orders cannot fall in either lattice quotient.

Write \(T=abcb\), \(U=cbab\), translations by \((-1,0)\), \((0,1)\).
The extra relators for \(H_n\) can be taken as \(T^{2^e},U^{2^e}\) in
the square case, and \(T^{2^e}U^{2^e},T^{2^e}U^{-2^e}\) in the diagonal
case. After deleting any marker, both \(T,U\) have order at most two:
deleting \(a\) leaves conjugates of \(c\), deleting \(c\) leaves
conjugates of \(a\), and deleting \(b\) leaves \(ac,ca\). Therefore
every lattice relator survives every deletion. The even Coxeter relators
also survive. Thus \(H_n\) has Boolean deletions and IP by UA.D.

Moreover,
\[
 \gamma_3(W)=2\mathbb Z^2,
 \qquad \gamma_3(H_n)=2\mathbb Z^2/\Lambda_n,
 \qquad |\gamma_3(H_n)|=2^{n-5}.                         \tag{S2}
\]
Indeed reduction modulo \(2\mathbb Z^2\) makes the adjacent squares
central; conversely \([[a,b],c]\) and \([[b,c],a]\) give the two doubled
unit translations, up to signs. Thus \(H_5\) is class two of order 32.
The squared translations are not always independent cyclic subgroups:
in \(H_6\) they coincide and have order two. Formula (S2) uses a lattice
index and does not require independence of those cyclic subgroups. \(\square\)

\section{Recognizing an affine presentation}\label{recognizing-an-affine-presentation}

We write an integral affine transformation as \((M,t)\), acting by
\(v\mapsto Mv+t\), with multiplication

\[ (M,t)(N,u)=(MN,t+Mu). \tag{1} \]

Let \(E\) be a group generated by finitely many such transformations, and
let \(Q\leq\operatorname{GL}_d(\mathbb Z)\) be its finite linear image.
Its translation kernel \(L\) is a subgroup of \(\mathbb Z^d\), invariant under
\(Q\). A translation is written \(t_v=(I,v)\).

\textbf{Lemma AL.1 (fixed Schreier lattice).} Choose a word lift \(\widehat q\in E\)
for each \(q\in Q\), with \(\widehat1=1\). Then \(L\) is generated by the
translations

\[\widehat{\overline s q}^{-1}s\widehat q,\qquad
s\text{ a distinguished generator},\quad q\in Q.\tag{2}\]

\emph{Proof.} The products have trivial linear image. Write a partial word
as \(\widehat q n\), with \(n\in L\). Left multiplication by \(s\) changes
this to \(\widehat{\overline s q}t n\), where \(t\) is the transition in
(2). Induction expresses a word of trivial linear image as a product
of these transitions; its final representative is \(\widehat1=1\).
Inverses introduce no extra generator: the transition for $s^{-1}$ at $q$ is the inverse of the transition for $s$ at $\overline{s}^{-1}q$. Thus the induction applies to words with positive and negative letters. Every transition itself is in \(L\), so they generate exactly \(L\).
\(\square\)

The normalization in (2) matters. The raw differences
\(s\widehat q\widehat{\overline s q}^{-1}\) need not generate \(L\)
without conjugation. Each is a conjugate of the corresponding transition
by \(\widehat{\overline s q}\). Consequently their \(Q\)-orbits generate
\(L\). In particular, if a proposed \(Q\)-invariant lattice contains every
raw difference, it contains every normalized transition and hence \(L\).
If exhibited raw differences also span the proposed lattice, equality
follows. This invariant-lattice version justifies the raw-difference
format of the unchanged earlier \(\{6,6\}\) checker {[}C66{]}. The new
\(\{6,4\}\) checker uses (2) directly.

This is a finite proof of a lattice identity. Containment in a proposed
lattice is checked on the normalized transitions, or on the raw differences
with the stated invariance condition. Equality follows when independent
integer combinations give a basis of that lattice, or when the greatest
common divisor of the full-rank minors gives its exact index. Sampling
finite quotients is not a replacement for this step.

\textbf{Lemma AL.2 (presentation bridge).} Let \(U\) be a finitely presented marked
group whose relators hold on the affine generators of \(E\). Suppose their
linear images define the same epimorphism to \(Q\). If an exact
Reidemeister--Schreier presentation proves
\(\ker(U\to Q)\cong\mathbb Z^d\), and the affine translation kernel \(L\)
has rank \(d\), then the marked epimorphism \(U\to E\) is an isomorphism.

\emph{Proof.} The restriction maps \(\ker(U\to Q)\) onto \(L\): an affine
translation has a preimage in \(U\), and its linear image is trivial.
This is a surjection between free abelian groups of rank \(d\). Tensoring
with \(\mathbb Q\) shows that its kernel has rank zero. A subgroup of a
free abelian group is torsion-free, so that kernel is trivial. The
remaining kernel of \(U\to E\) lies in this subgroup, and is therefore
trivial as well. \(\square\)

The hypothesis concerns the actual kernel presentation. Computing only
its abelianization would not establish it: a nonabelian group can have
free abelian abelianization of rank \(d\). Our applications retain the
exact rewriting and Tietze algorithms used to verify this finite
presentation step. It is a computer-assisted input, not a Lean proof.

\section{Periods and intersections of distinguished subgroups}\label{periods-and-intersections-of-distinguished-subgroups}

Let \(w_1,\ldots,w_d\in U=E\) be translation words with linearly independent
vectors \(v_1,\ldots,v_d\). Write
\(\Lambda=\sum_i\mathbb Z v_i\leq L\) and suppose it is \(Q\)-invariant.
For a positive integer \(m\), adjoining \(w_i^m=1\) to the complete
presentation of \(U\) gives exactly

\[ E_m=E/(m\Lambda).\tag{3}\]

Indeed, the normal closure of the powers is \(m\Lambda\): conjugation by
an affine element acts on translations by its linear part, and this
lattice is already invariant. Hence

\[ |E_m|=|Q|[L:\Lambda]m^d.\tag{4}\]

This proves equality for the presented group, not merely the order of a
chosen quotient. If \(Q\) is solvable, every \(E_m\) is solvable.

Suppose now that the marking has rank three with distinct generators,
its endpoints commute, both adjacent products have order at least three, and
its two adjacent subgroups \(A=E_{01}\) and \(B=E_{12}\) are finite and
inject into \(Q\). For each \(q\in\overline A\cap\overline B\), let
\((q,t_A(q))\) and \((q,t_B(q))\) be their unique affine lifts. Put

\[ \Delta(q)=t_A(q)-t_B(q).\tag{5} \]

\textbf{Theorem AL.3 (period-lattice intersection criterion).} The images of \(A\)
and \(B\) in \(E_m\) meet precisely in the elements indexed by

\[ \{q\in\overline A\cap\overline B:\Delta(q)\in m\Lambda\}.\tag{6}\]

In particular, if \(A\cap B=\langle r_1\rangle\) and every other value
\(\Delta(q)\) lies outside \(\Lambda\), then every \(E_m\), \(m\geq1\), has
the full rank-three intersection property and the same adjacent orders
as \(E\). Its order is (4).

\emph{Proof.} Equality of two cosets modulo translations first forces equality
of their linear parts. For equal linear part, their affine quotient is
translation by the difference (5). This proves (6). Each adjacent
subgroup remains injective in the quotient because its linear action is
faithful, so its order and the product order of its two reflections are
unchanged. The assumed exclusions apply to \(m\Lambda\subseteq\Lambda\).
Thus the adjacent intersection is exactly the common involution subgroup,
and the rank-three test proves all subset intersections. \(\square\)

The proof separates two possibilities often hidden by equal group-order
counts. The point-group intersection may already have order two, as in
the \(\{6,4\}\) application. Or it may be larger, with a short list of
nonzero translation differences, as in the \(\{6,6\}\) application. Those
differences must be excluded from the period lattice explicitly.

\subsection*{One coarse quotient controls all smaller periods}
The same argument has a group-theoretic formulation which is useful independently of coordinates. Let $A,B,L$ be subgroups of a group $E$, with
\[
 L\cap A=L\cap B=1.
\]
Let $N\unlhd E$ lie in $L$. If the images of $A,B$ in $E/N$ intersect exactly in the image of $A\cap B$, then the same statement holds in $E/M$ for every normal subgroup $M\le N$.

Indeed, suppose $a\in A$, $b\in B$ have the same image modulo $M$. They have the same image modulo $N$, so there is $c\in A\cap B$ with $aN=cN=bN$. Then $ac^{-1}\in N\cap A=1$ and $bc^{-1}\in N\cap B=1$, proving $a=b=c$. The reverse inclusion is immediate. In the affine setting $L$ is the translation kernel, $N=\Lambda$ and $M=m\Lambda$. Point-faithfulness supplies precisely the two trivial intersections. This statement, for an arbitrary indexed family of smaller normal periods, is verified in \texttt{AffinePeriodIntersections.lean}. The finite coordinate checks which identify $L,\Lambda$ and the initial overlap remain separate mathematical obligations.

\section{Changing one adjacent order}\label{changing-one-adjacent-order}

Let \((B;a,b,c)\) be a rank-three string C-group of type \(\{h,q\}\), with
\(h,q\) even, and suppose the map

\[\rho:B\longrightarrow B_{12},\qquad a\mapsto1, b\mapsto b, c\mapsto c
\tag{7}\]

is a retraction. For \(k\) a positive multiple of \(q\), take the standard
marked epimorphism \(D_{2k}\to B_{12}=D_{2q}\). The compatible product

\[F=\{(g,d)\in B\times D_{2k}:\rho(g)=\overline d\}\tag{8}\]

has markers \((a,1),(b,\widetilde b),(c,\widetilde c)\). The evenness of
\(h\) is forced by (7); it also ensures that the first pair has unchanged
order \(h\).

\textbf{Theorem AL.4 (one-edge extension).} These markers generate \(F\). If \(B\) is finite, its order
is \(|B|k/q\). Its type is \(\{h,k\}\), and it has the intersection
property. The projection onto \(B\) is injective on the first adjacent
subgroup. If \(B\) is solvable, so is \(F\).

\emph{Proof.} The word \((r_1r_2)^q\) generates the entire cyclic kernel of
(8), since its first coordinate is trivial and its second has order
\(k/q\). Surjectivity onto \(B\) then proves generation and the order formula.
The first adjacent subgroup has second coordinate of order at most two,
with that coordinate determined by the image under (7); thus its
projection is injective. The second adjacent subgroup is the full
\(D_{2k}\). If an element belongs to both, its projection belongs to
\(\langle b\rangle\) by IP in \(B\). Subtract the appropriate common
involution. The remaining element lies in the first subgroup and in the
projection kernel, hence is trivial. This proves IP by the rank-three
criterion. Solvability follows from the cyclic kernel. \(\square\)

Theorem AL.4 is a specialization of UA.SK; it is included to make clear that
only the indicated deletion is needed here. More simultaneous edges use
the supported coordinate kernels of that theorem.

There is a useful exact presentation variant. Suppose a complete marked
presentation for \(B\) contains \((bc)^q=1\). Replace that relation by
\((bc)^k=1\), retain the remaining relators, and add
\([a,(bc)^q]=1\). Call the resulting presented group \(H\). If every retained
relator holds in (8), then \(H=F\) as a marked group. Indeed,
\(C=\langle(bc)^q\rangle\) is normal in \(H\): the last relation centralizes
it by \(a\), and \(b,c\) invert it. Its order is at most \(k/q\), and \(H/C=B\)
by the presentation. Consequently \(|H|\leq|B|k/q\). The marked
surjection onto (8) attains this bound, proving equality and injectivity.
The same argument is valid for an infinite base if the finite kernel
maps injectively and the induced map on the base is an isomorphism.

\section*{Attribution and verification}
The affine lattice method, Reidemeister--Schreier rewriting and Tietze transformations are classical. Hou--Feng--Leemans already used free abelian normal subgroups and their power subgroups. Our earlier cubic calculation and supported-kernel theory are incorporated here with their full proofs and exact computational dependencies. The companion application identifies every source word and exceptional parameter.

The regular-polytope correspondence is retained from McMullen--Schulte, \emph{Abstract Regular Polytopes} (2002), Section 2E. The local rank-three intersection test is proved here. The square affine presentation has its elementary normal-form proof here; the cubic presentations require the two complete exact GAP kernel programs. No computation of an abelianization is substituted for a presentation of the actual kernel.

The citation [C66] denotes the project's earlier \emph{One affine model for all three $\{6,6\}$ lattice families}; its complete local proof and two programs are incorporated into the companion application. No additional statement from that note is assumed without its proof here.

The full common paper, the application, and all six mathematical programs belong to the replacement proof. A theorem already available in another topic is not free in an independent comparison. The source's historical programs have not been supplied, so their cost is unknown. No compression factor is asserted in this edition.
\end{document}
